%! Equations and Graphs in Both Directions — Unit 2, Lesson 2.3 — Group Activity (Answer Key)
\documentclass[10pt]{article}
\usepackage{atda-article}
\usepackage{atda-key}   % pulls in -boxes; do NOT also load -boxes
\newcommand{\TallMath}[1]{$\displaystyle #1\rule[-1.4em]{0pt}{3.2em}$}
\setlength{\workrowsep}{6pt}

\begin{document}

\pageheader{Unit 2, Lesson 2.3}{Group Activity --- Answer Key}
% NAMESTRIP: no \namedateperiod here — mirrors the blank. Teacher-only prose goes
% in the lesson plan as \begin{teachernote}[Group Activity], never in a key.

\vspace{0.15in}

\begin{scenariobox}[The Scenario]{royal}
\small
Half of today's work goes equation $\to$ graph, and half goes graph $\to$ equation. Every time you
write an equation, \textbf{finish with a substitution check}: pick one point you can read off the
graph, put its $x$ into your equation, and confirm you get its $y$. An unchecked equation is a
guess. \textbf{Everyone starts at Tier R.} When your whole group can do Tier R without help, move on.
\end{scenariobox}

\vspace{0.10in}

\boxguard[30]
\begin{tcolorbox}[colback=white, colframe=black!40, title=\textbf{Tier R --- Remediate},
    fonttitle=\bfseries, arc=2mm, left=3mm, right=3mm, top=2mm, bottom=2mm, breakable]
\small
\textbf{Equation $\to$ graph.} Build the graph of $y=(x+1)^2-4$ from the parent.

\begin{enumerate}[leftmargin=1.6em, itemsep=4pt, topsep=2pt]
  \item Name the parent, then give the anchor point of the image. Say which slide came from the
        number inside the parentheses and which came from the number outside.

\ansline{Parent $f(x)=x^2$; anchor point $(-1,-4)$. The $+1$ is \textbf{inside}, so it slid the graph}\\
\ansline{\textbf{left $1$}; the $-4$ is \textbf{outside}, so it slid the graph \textbf{down $4$}.}\\

  \item Complete the table, then plot your five points on the grid and join them with a smooth
        curve. The parent is already drawn.
\smallskip

\renewcommand{\arraystretch}{1.3}
\begin{tabular}{@{}|c|c|c|c|c|c|@{}}
\hline
$x$ & $-3$ & $-2$ & $-1$ & $0$ & $1$ \\ \hline
$y=(x+1)^2-4$ & \ans{$0$} & \ans{$-3$} & \ans{$-4$} & \ans{$-3$} & \ans{$0$} \\
\hline
\end{tabular}
\smallskip

\begin{center}
\begin{tikzpicture}
\begin{axis}[
    width=9.4cm, height=5.4cm,
    xlabel={$x$}, ylabel={$y$},
    xmin=-4.4, xmax=3.4, ymin=-5.4, ymax=5.4,
    xtick={-4,-3,-2,-1,0,1,2,3}, ytick={-5,-4,-3,-2,-1,0,1,2,3,4,5},
    grid=both, grid style={linegray!45}, axis lines=left,
    tick label style={font=\tiny}, label style={font=\scriptsize},
    legend entries={{$y=x^2$ (the parent)},{$y=(x+1)^2-4$}},
    legend style={font=\scriptsize, draw=none, fill=none,
      at={(0.5,1.02)}, anchor=south, legend columns=2}]
  \addplot[royal, very thick, domain=-2.2:2.2, samples=80]{x^2};
  \addplot[keyred, very thick, dashed, domain=-3.9:1.9, samples=80]{(x+1)^2-4};
  \addplot[keyred, only marks, mark=*, mark size=1.6pt] coordinates
    {(-3,0) (-2,-3) (-1,-4) (0,-3) (1,0)};
\end{axis}
\end{tikzpicture}
\end{center}

  \item \textbf{Graph $\to$ equation, the easy direction.} A parabola has exactly the same shape as
        $y=x^2$, and its lowest point is at $(4,0)$. Write its equation, then check it by
        substituting $x=6$.

\ansline{Same shape means $a=1$; the vertex $(4,0)$ gives $h=4$ and $k=0$, so $y=(x-4)^2$.}\\
\ansline{Check: at $x=6$, $y=(6-4)^2=4$, so the graph must pass through $(6,4)$ --- it does.}\\

  \item Which of the domain and the range did the $-4$ in $y=(x+1)^2-4$ move? Give the range of the
        image in interval notation.

\ansline{The $-4$ is outside, so it moved the \textbf{range}: $[0,\infty) \to [-4,\infty)$. Domain stays $(-\infty,\infty)$.}\\
\end{enumerate}
\end{tcolorbox}

\vspace{0.10in}

\boxguard[30]
\begin{tcolorbox}[colback=white, colframe=black!40,
    title=\textbf{Tier A --- Approaching Proficiency},
    fonttitle=\bfseries, arc=2mm, left=3mm, right=3mm, top=2mm, bottom=2mm, breakable]
\small
\textbf{Graph $\to$ equation.} Now the pictures are handed to you and the shapes are not all the
parent's.

\begin{center}
\begin{tabular}{@{}c@{\hspace{0.6cm}}c@{}}
\begin{tikzpicture}
\begin{axis}[
    width=6.6cm, height=4.4cm, title={\scriptsize (a) a parabola},
    xmin=0, xmax=6.4, ymin=-4.4, ymax=11.4,
    xtick={0,1,2,3,4,5,6}, ytick={-3,0,3,6,9}, minor y tick num=2,
    grid=both, grid style={linegray!45}, minor grid style={linegray!30},
    axis lines=left,
    tick label style={font=\tiny}]
  \addplot[royal, very thick, domain=1:5, samples=90]{3*(x-3)^2-2};
\end{axis}
\end{tikzpicture}
&
\begin{tikzpicture}
\begin{axis}[
    width=6.6cm, height=4.4cm, title={\scriptsize (b) an exponential},
    xmin=-3.4, xmax=2.4, ymin=0, ymax=7.4,
    xtick={-3,-2,-1,0,1,2}, ytick={0,2,4,6}, minor y tick num=1,
    grid=both, grid style={linegray!45}, minor grid style={linegray!30},
    axis lines=left,
    tick label style={font=\tiny}]
  \addplot[linegray, thick, dashed, domain=-3.3:2.3, samples=2]{2};
  \addplot[royal, very thick, domain=-3.3:2.2, samples=90]{2^x+2};
\end{axis}
\end{tikzpicture}
\end{tabular}
\end{center}

\begin{enumerate}[leftmargin=1.6em, itemsep=4pt, topsep=2pt]
  \item \textbf{Graph (a).} Read the vertex, then use the point $(4,1)$ to find $a$. Write the
        equation in transformation form.
\begin{work}
  1 &= a(4-3)^2-2 \\
  1 &= a-2 \\
  3 &= a
\end{work}

  \item Check your equation from item 1 at $x=5$, and say what the graph must show there if you are
        right.

\ansline{Vertex $(3,-2)$, so $y=3(x-3)^2-2$. At $x=5$: $3(5-3)^2-2=3(4)-2=10$.}\\
\ansline{The graph must pass through $(5,10)$ --- and it does, so the equation is verified.}\\

  \item \textbf{Graph (b).} The dashed line is the flat line this curve crawls along. Give its
        equation, give the point the curve passes through on the $y$-axis, and write the function.
        Then state the range in interval notation --- bracket or parenthesis, and why.

\ansline{The flat line (asymptote) is $y=2$; the curve crosses the $y$-axis at $(0,3)$, since $2^0+2=3$.}\\
\ansline{So $y=2^x+2$: the parent $2^x$ translated \textbf{up $2$}.}\\
\ansline{Range $(2,\infty)$ --- a \emph{parenthesis}, because $2^x$ never equals $0$, so $y$ never equals $2$.}\\

  \item \textbf{In context.} A gym charges a joining fee plus a flat monthly rate. The graph of the
        total cost $C(x)$ after $x$ months is a straight line through $(0,40)$ and $(4,100)$.
        \begin{enumerate}[label=(\alph*), leftmargin=1.6em, itemsep=3pt, topsep=3pt]
          \item Find the monthly rate from the two points, and write $C(x)$.

\ansline{Rate $=(100-40)/(4-0)=\$15$ per month; joining fee $\$40$. So $C(x)=15x+40$.}\\

          \item Use your equation to find the total cost after $7$ months.
\begin{work}
  C(7) &= 15(7)+40 \\
       &= 105+40 \\
       &= 145
\end{work}

          \item The notes said that for a \emph{line} you report two points rather than naming a
                transformation. Give the reason, in your own words.

\ansline{Different transformations of $y=x$ give the identical line: $2f(x)$ and $f(2x)$ are both}\\
\ansline{$y=2x$, and sliding a line sideways looks the same as sliding it up or down. Two points are unambiguous.}\\
        \end{enumerate}
\end{enumerate}
\end{tcolorbox}

\vspace{0.10in}

\boxguard[30]
\begin{tcolorbox}[colback=white, colframe=black!40, title=\textbf{Tier E --- Extension},
    fonttitle=\bfseries, arc=2mm, left=3mm, right=3mm, top=2mm, bottom=2mm, breakable]
\small

\begin{enumerate}[leftmargin=1.6em, itemsep=4pt, topsep=2pt]
  \item \textbf{Find the typo.} Each student below typed something into a graphing calculator and
        got a graph that does not match what they were asked for. Name the error and give the
        keystrokes that fix it.
        \begin{enumerate}[label=(\alph*), leftmargin=1.6em, itemsep=3pt, topsep=3pt]
          \item Asked for: the exponential parent moved \emph{right} $5$. Typed:
                \texttt{Y1=2\^{}X-5}.

\ansline{The $5$ landed \emph{outside} the exponent, so the screen shows $2^x$ moved \textbf{down $5$}, not right.}\\
\ansline{Fix: \texttt{Y1=2\^{}(X-5)}. Check at $x=5$: $2^{5-5}=1$, matching the parent's $2^0=1$.}\\

          \item Asked for: the quadratic parent \emph{reflected over the $x$-axis}. Typed:
                \texttt{Y1=(-X)\^{}2}. The screen showed a parabola opening \emph{upward}.

\ansline{\texttt{(-X)\^{}2} negates \emph{first}, then squares, and $(-x)^2=x^2$ --- so it drew the parent again.}\\
\ansline{Fix: \texttt{Y1=-X\^{}2} (using the \texttt{(-)} key), which squares first and then negates.}\\
        \end{enumerate}

  \item \textbf{The window.} A classmate typed $y=(x-15)^2-40$ correctly and pressed
        \texttt{ZOOM}\,\texttt{6}, which sets the standard window $-10 \le x \le 10$,
        $-10 \le y \le 10$. The screen came up empty.
        \begin{enumerate}[label=(\alph*), leftmargin=1.6em, itemsep=3pt, topsep=3pt]
          \item Explain why the screen is empty, using the anchor point.

\ansline{The anchor point is $(15,-40)$ --- outside the window on both axes, so the whole curve misses the screen.}\\

          \item Give four window values that would show the vertex and both sides of the curve, and
                say how you chose them \emph{from the equation}, without graphing anything.

\ansline{The equation gives $h=15$ and $k=-40$, so bracket them: \texttt{Xmin}$=5$, \texttt{Xmax}$=25$,}\\
\ansline{\texttt{Ymin}$=-50$, \texttt{Ymax}$=70$ (since $y=60$ at $x=5$ and at $x=25$, that shows both arms).}\\
        \end{enumerate}

  \item \textbf{Justify.} A classmate says: ``My calculator screen looks exactly like the printed
        graph, so my equation must be right.'' Explain why that is not proof --- use $y=(x-3)^2$ and
        $y=(x-3)^2+0.5$ as your example --- and then say what \emph{would} be proof.

\ansline{Those two graphs sit $0.5$ apart everywhere. On a screen showing $-10$ to $10$ vertically in about}\\
\ansline{$60$ rows of pixels, half a unit is a pixel or two --- so the curves look identical yet are different}\\
\ansline{functions. Proof is a substitution: at $x=3$ one gives $0$, the other $0.5$.}\\
\end{enumerate}
\end{tcolorbox}

\end{document}
